\section{Acquired geometry, singular mass and state-dependent waiting}
\label{sec:geometry}
The following quantization calculation is inherited from the stopped-experiment analysis \cite{GTF52,GTF53}; we include its proof, constants' range and singular extension to keep the present transfer self-contained. Put
\[
 \sigma=(f,\sqrt w)_\#\nu.
\]
Assume its location support is a compact $K\subset H$ and $0\le\sqrt w\le1$. Fix $d,\alpha>0$, $r_0\in(0,1]$, and $c,C>0$ such that for every $x\in K$, $0<r\le r_0$ and $0<s\le1$,
\begin{equation}\label{eq:rectangle}
 cr^d s^\alpha\le\sigma\bigl(B(x,r)\times[0,s]\bigr)
                  \le Cr^d s^\alpha.
\end{equation}
Thus the geometry is that of the actual marked law. In particular it can be singular in its ambient Euclidean space. A fixed finite cover of $K$ by radius $r_0$ balls gives $\sigma(H\times[0,s])\le C_Ts^\alpha$.

\begin{proposition}[Marked selection and quantization]\label{prop:geometry}
Under \eqref{eq:rectangle}, uniformly for $b\ge1$, $m\ge1$,
\begin{equation}\label{eq:regular-profile}
 \Qm_m(b)\asymp
 F(b,m):=\max\{b^{-2/\alpha},(bm)^{-2/(d+\alpha)}\}.
\end{equation}
Constants depend only on the displayed geometric data and the diameter of $H$, not on $b,m$.
\end{proposition}
\begin{proof}
A mass-one density at most $b$ cannot put more than $bC_Ts^\alpha$ mass at reliability at most $s$. Taking $s=(2C_Tb)^{-1/\alpha}$, reduced by a fixed constant if necessary, gives a lower of order $b^{-2/\alpha}$.

For any dictionary of $m$ centers in $H$, a marked distortion sublevel $d_C\le t^2$ lies in the union of $m$ rectangles with reliability at most $t$ and location radius $t$. A center not in $K$ can be moved for this estimate to a point of $K$ within $t$, expanding the radius to $2t$. For $2t\le r_0$ its mass is at most $C_Rm t^{d+\alpha}$. Set $t=c_0(bm)^{-1/(d+\alpha)}$, with
\[
 c_0\le \min\{r_0/2,1,(2C_R)^{-1/(d+\alpha)}\}.
\]
At least half the selected mass has distortion greater than $t^2$. This proves the other lower for the entire range $b,m\ge1$.

The marginal at $s=1$ has two-sided ball mass of order $r^d$. A maximal separated-set argument supplies, for every integer $k\ge1$, $k$ points in $K$ separated by at least $a k^{-1/d}$, with one fixed $a>0$. For completeness, a maximal $r$-separated set covers $K$ by $r$ balls, so its cardinality is at least $(Cr^d)^{-1}$ for $r\le r_0$; choose $r$ a sufficiently small multiple of $k^{-1/d}$ and retain $k$ points. The $k=1$ case imposes no separation condition.

Choose $C_1$ so that $cC_1^{d+\alpha}\ge1$, then choose $c_1\in(0,1]$ so small that
\[
 C_1 c_1^{\alpha/[d(d+\alpha)]}\le a/3.
\]
Increase the explicit threshold
\begin{equation}\label{eq:b-threshold}
 b_0=\max\left\{1,(2/c_1)^{\alpha/d},
        (C_1/\min\{r_0,1\})^{d+\alpha}\right\}
\end{equation}
if needed to accommodate the chosen packing constant. For $b\ge b_0$ set
\[
 k=\min\{m,\lfloor c_1 b^{d/\alpha}\rfloor\},\qquad
 t=C_1(bk)^{-1/(d+\alpha)}.
\]
The radius satisfies $t\le\min(r_0,1)$ and $3t\le a k^{-1/d}$. The $k$ reliability rectangles are therefore disjoint and their union has mass at least $kc t^{d+\alpha}\ge1/b$. On this union the dictionary formed by the $k$ centers has distortion at most $2t^2$. Normalize its restriction to mass one; its density is at most $b$. If $k=m$ this is the second term of $F$; otherwise $k$ is bounded below by a fixed multiple of $b^{d/\alpha}$ and it is the first term.

For $1\le b<b_0$, use density one and a single center. Its distortion is at most $1+\operatorname{diam}(H)^2$, while $F(b,m)\ge b_0^{-2/\alpha}$. This handles every $m$, including arbitrarily large $m$; it is a bounded-budget argument, not an assertion that the remaining $(b,m)$ set is compact.
\end{proof}

\begin{corollary}[Matching curves under regular acquired geometry]\label{cor:regular}
Assume Theorem~\ref{thm:main}, \eqref{eq:rectangle}, and $\kappa_+/\kappa_-\le\Gamma<\infty$. Then, for $N\ge\kappa_+$,
\[
 R_\cp^\Theta(N,m)-B_0\asymp F(N/\kappa_+,m).
\]
If $B$ mandatory physical contexts verify Proposition~\ref{prop:semantic} and the upper uses $B$ live labels, then for every $M\ge B+1$,
\begin{equation}\label{eq:matched-aut}
 R_\aut^\Theta(N,M)-B_0\asymp F(N/\kappa_+,M-B).
\end{equation}
The constants additionally depend on $\ell,u,\Gamma,k_-,k_+$, but are independent of $N,M$ and of the absolute scale of $\kappa_+$.
\end{corollary}
\begin{proof}
Apply Proposition~\ref{prop:geometry} to the two arguments of Theorem~\ref{thm:main}. Multiplying $b\ge1$ by fixed factors changes either power by fixed factors. The upper endpoint $\max(1,\ell b)$ obeys the same comparison. Mandatory width gives the same $M-B$ on both sides. No regular variation in an ambient parameter dimension is used.
\end{proof}

\begin{theorem}[An actual atomic transition]\label{thm:atomic}
Let the marked acquired law be
\[
 \nu=\lambda\delta_{(x_*,0)}+(1-\lambda)\nu_0,
 \qquad 0<\lambda<1,
\]
where $\nu_0$ is a probability law satisfying \eqref{eq:rectangle}; the notation here identifies marks with their location and reliability. Write $r=(1-b\lambda)_+$. If $r=0$, then $\Qm_m^\nu(b)=0$ for $m\ge1$. If $r>0$ and $m\ge2$, then
\begin{equation}\label{eq:atomic-profile}
 r\Qm_m^{\nu_0}\!\left(\frac{b(1-\lambda)}r\right)
 \le\Qm_m^\nu(b)
 \le r\Qm_{m-1}^{\nu_0}\!\left(\frac{b(1-\lambda)}r\right).
\end{equation}
Together with \eqref{eq:regular-profile}, this is a matching curve, including its vanishing at the actual-mass threshold $b\lambda=1$.
\end{theorem}
\begin{proof}
The atomic selected mass is at most $b\lambda$, so the regular part has mass at least $r$. Removing surplus regular mass can only decrease its nonnegative distortion. A regular measure of mass exactly $r$, normalized to probability, has density at most $b(1-\lambda)/r$ relative to $\nu_0$. This proves the lower with the original dictionary. For the upper put selected mass $b\lambda$ on the atom, use one center at $x_*$, and use the other $m-1$ centers for an optimal regular submeasure of mass $r$. If $b\lambda\ge1$, select the atom alone with density $1/\lambda\le b$. For $m\ge2$, $m-1$ and $m$ differ by at most a factor two, proving matching through Proposition~\ref{prop:geometry}. The factor $r$ and the threshold $b\lambda$ cannot be removed by renormalizing the regular stratum.
\end{proof}

For example, a product of an Ahlfors-regular Cantor location law and the reliability law $\Pp(T\le s)=s^\alpha$ satisfies \eqref{eq:rectangle}. Products with unequal factor dimensions give anisotropic cell constructions before the scalar distortion is chosen. An actual revealing atom can be supplied as an additional reported outcome. These are singular acquired laws, not smooth parameter-space surrogates. Under conditional density envelopes the two sides of Theorem~\ref{thm:main} retain their respective atomic thresholds; without additional exact-law structure, they must not be collapsed into one threshold.

\begin{theorem}[Unknown, history-dependent heralding]\label{thm:herald}
There are two physical phases, gate and acquisition. A gate call costs one and permits no stopping. Its conditional success probability $p_h$, allowed to depend on the entire history, command and unknown fixed model, satisfies
\[
 0<p_-\le p_h\le p_+\le1.
\]
Failure leaves the gate phase; success enables one acquisition call. The latter costs one, is productive, and returns to the gate phase on rejection. Only the corresponding physical command is legal at each of the two ports; both ports are mandatory under one admitted model. Assume the task comparison and conditional mark envelopes of Section~\ref{sec:setup}. Then Theorem~\ref{thm:main} applies with
\begin{equation}\label{eq:kappa}
 B=2,\qquad \kappa_-=1+1/p_+,\qquad \kappa_+=1+1/p_-.
\end{equation}
In particular \eqref{eq:matched-aut} holds under \eqref{eq:rectangle} and $p_+/p_-\le\Gamma$. The program needs only the certified bounds, not the true $p_h$. The hidden system is not reset between attempts.

No procedure is feasible for $N<\kappa_-$. If the family contains a model with gate success constantly $p_-$ and the same mandatory protocol, robust feasibility is exactly $N\ge\kappa_+$. Uniformity as $p_-\downarrow0$ is on this rescaled feasible budget. At fixed finite $N$, that family's problem becomes infeasible when $1+1/p_->N$. At $p_-=p_+=1$ both one gate and one acquisition call are paid, so $\kappa_+=2$.
\end{theorem}
\begin{proof}
For the all-strategy effort lower, put $V=0$ in gate phase and $V=-1/p_+$ in acquisition phase. After every productive call, committed or not, give it value zero. On a gate call,
\[
 \E[1+V_{\rm next}-V_{\rm current}\mid h]=1-p_h/p_+\ge0.
\]
On an acquisition call the increment is $1+1/p_+=\kappa_-$. These are precisely \eqref{eq:effort-drift}, without a fixed transition row.

For an attempt program, repeat the gate command until success and then acquire. The two-label cost supersolution is $W_G=\kappa_+$, $W_A=1$. Indeed
\[
 1+(1-p_h)W_G+p_hW_A\le W_G,\qquad 1\le W_A.
\]
Lemma~\ref{lem:navigator} proves the uniform conditional attempt bound. Proposition~\ref{prop:semantic} gives two live labels in the lower, since these are different physical ports, not aliases. Thus all hypotheses of Theorem~\ref{thm:main} have been checked from conditional raw laws. When the constant-$p_-$ model belongs to the family, predictable gate counts give expected cost at least $\kappa_+$ for one productive call; accepting that first call attains it. This proves the exact feasibility statement. No independent-cycle or true-rate argument appears.
\end{proof}
